\origpage{42}{ROBERTSON AND SEYMOUR}
\noindent values, with $i<i'$ say, then $i<T$ and $i'>0$, and we have that $Y_{i+1}$ does not
separate $v$ and $\rho$, and $Y_{i'}$ separates $v$ and $\rho$. Let $P$ be a path in $G$ from $v$ to $\rho$
avoiding $Y_{i+1}$. Then $P$ meets $Y_{i'}$, and so $Y_{i+1}$ does not separate $Y_{i'}$ and $\rho$,
contrary to our hypothesis.

It follows that for every vertex $v$ there is a unique value of $i$ with $v\in W_i$,
and so $W_0, W_1,..., W_T$ partition $V(G)$. Clearly, $Y_i\subseteq W_i$ $(1\le i\le T)$. It
remains to prove that the sequence
\begin{equation*}
W_0\cup Y_1, W_1\cup Y_2,..., W_{T-1}\cup Y_T, W_T
\end{equation*}
is a path-decomposition of $G$. We need only show (W1), since (W2) is
obvious.

Let $e$ be an edge of $G$, with ends $u, v$ say. Let $u\in W_i$, $v\in W_j$, where $i\le j$
say. If $i=j$, then \del{$\rho$}\fix{$e$}\fixnote{fix:1.7-proof}{The edge $e$ is meant.} has both ends in some set in the sequence, as required;
we assume therefore that $i<j$. Thus $Y_{i+1}$ separates $v$ and $\rho$. It does not
separate $u$ and $\rho$, since $u\in W_i$, and yet $u, v$ are adjacent. It follows that
$v\in Y_{i+1}$, and so $e$ has both ends in $W_i\cup Y_{i+1}$, as required.

\sectionhead{2. \textsc{Grids}}

If $\theta\ge 2$, the \textit{$\theta$-grid} is the simple graph with vertices $v_{ij}$ $(1\le i, j\le\theta)$ in
which $v_{ij}$ and $v_{i'j'}$ are adjacent if $|i-i'|+|j-j'|=1$. $\mathscr{F}_\theta$ denotes the set of
all graphs with no minor isomorphic to the $\theta$-grid. For the remainder of the
paper, $\theta\ge 6$ will be fixed and, for convenience, even. In this section we prove
some connectivity results about the graphs in $\mathscr{F}_\theta$.
\stmt
(2.1). \textit{If $P_1,..., P_m$ are disjoint paths of a graph $G\in\mathscr{F}_\theta$, and $B_1,..., B_n$ are
also disjoint paths, and each $P_i$ has exactly one vertex in common with each
$B_j$, and each $P_i$ meets $B_1,..., B_n$ in that order, then either $m<\theta$ or
$n<\frac12\theta\cdot\theta!$.}
\stmt
\textit{Proof}. We assume that $m\ge\theta$. Since there are only $\frac12(\theta!)$ different
orderings of $P_1,..., P_\theta$ up to reversal, it follows that there exists $J\subseteq\{1,..., n\}$
with $|J|\ge n/\frac12(\theta!)$ such that for $j, j'\in J$, $B_j$ and $B_{j'}$ meet $P_1,..., P_\theta$ in the
same order or the reverse. If $n\ge\frac12\theta\cdot\theta!$, then $|J|\ge\theta$ and so $G$ has a minor
isomorphic to the $\theta$-grid, which is impossible. Thus $n<\frac12\theta\cdot\theta!$.
\stmt
(2.2). \textit{If $P_1,..., P_m$ are disjoint paths of $G\in\mathscr{F}_\theta$, and $B_1,..., B_n$ are disjoint
stars, and each $P_i$ has exactly one vertex in common with each $B_j$, and no $P_i$
contains the centre of any star $B_j$, and each $P_i$ meets $B_1,..., B_n$ in that order,
then either $m<\theta$ or $n<\theta(\theta-1)$.}
\stmt
(A \textit{star} is a tree in which one vertex, called the \textit{centre}, is adjacent to all
others.)
